\documentclass[11pt,a4paper]{article}

% --------- Packages ----------
\usepackage[left=1in,right=1in,top=0.75in,bottom=0.75in]{geometry}
\usepackage{titlesec}
\usepackage{enumitem}
\usepackage{hyperref}
\usepackage{setspace}
\usepackage{amsmath}

% --------- Formatting ----------
\setlength{\parindent}{0pt}
\setlist[itemize]{noitemsep, topsep=2pt, leftmargin=*}

\titleformat{\section}
{\large\bfseries}
{}
{0em}
{}
[\titlerule]

\begin{document}

% --------- Header ----------
\begin{center}
{\LARGE \textbf{SASWATA DAS}} \\[4pt]
Email: \href{mailto:saswata.das@students.iiserpune.ac.in}
{saswata.das@students.iiserpune.ac.in}
\quad | \quad
Phone: +91-8617250323
\end{center}

\vspace{4pt}

% --------- Profile ----------
\section*{Profile}
 
I am a 3rd-year BS-MS Physics student at IISER Pune interested in observational and computational cosmology. I have used Python to select stellar-stream members from Gaia DR3, analyze IllustrisTNG simulation data, and implement MCMC samplers, and I want to apply statistical inference to dark-matter and cosmological constraints.
 
% --------- Research Interests ----------
\section*{Research Interests}
 
Observational and computational cosmology, dark-matter substructure, stellar streams and Galactic dynamics, astronomical survey data analysis, CMB data analysis, cosmological parameter inference, and statistical methods for astrophysics.
 
% --------- Education ----------
\section*{Education}
 
\textbf{BS-MS, Physics} \hfill 2024 -- 2029 \\
Indian Institute of Science Education and Research (IISER) Pune \\
CGPA: 8.9 / 10
 
% --------- Research Projects ----------
\section*{Research Projects}
 
\textbf{KSP 7.0 -- Ripples in the Dark: Characterizing Stellar Streams as Dark Matter Substructure Probes}
\hfill Jun -- Jul 2026 \\
Indian Institute of Technology Bombay (IITB) \\
Mentors: Tamojeet Rowchoudhury (UC Berkeley), Ravi Kumar (University of Minnesota); Coordinator: Vihang Vidwans (IITB) \\
Code: \href{https://github.com/Saswata-Das117/KSP-7.0-Stellar-Stream-Dark-Matter}{github.com/Saswata-Das117/KSP-7.0-Stellar-Stream-Dark-Matter}

\begin{itemize}
\item Built a \textbf{GD-1 stellar-stream member-selection pipeline} on \textbf{Gaia DR3}, combining astrometric, photometric, proper-motion, and theoretical isochrone cuts to isolate candidate stream stars.
\item Transformed Gaia positions and proper motions into stream-aligned coordinates and characterized the proper-motion distribution of GD-1.
\item Measured the stream's density structure using \textbf{kernel density estimation (KDE)} and stream-track fitting to search for gaps and underdensities that could signal dark-matter substructure.
\item Analyzed \textbf{IllustrisTNG TNG50-1} simulation data to identify Milky Way-like host halos and quantify the dark-matter subhalos capable of perturbing streams.
\item Reduced \textbf{3.01 million} Gaia DR3 sources to \textbf{722 candidate} GD-1 stream members using photometric, proper-motion, spatial-footprint, and density cuts; identified \textbf{two underdensities} at \textbf{$\boldsymbol{\phi_1\approx-42^\circ}$} and \textbf{$\boldsymbol{-21^\circ}$} with significances of \textbf{2.7$\sigma$} and \textbf{4.0$\sigma$}.
\item Propagated Gaia proper-motion uncertainties through the full selection with a \textbf{Monte Carlo} (100 realizations); recovered the \textbf{$\boldsymbol{-21^\circ}$} gap in \textbf{99\%} and the \textbf{$\boldsymbol{-42^\circ}$} gap in \textbf{93\%} of realizations, with membership varying by $\sim$$^{+35}_{-34}$ stars around a median of \textbf{643}.
\item Extended the pipeline to \textbf{six other streams} (Indus, C-12, Ylgr, Orphan-Chenab, NGC~6397, Jhelum) with iterative sigma-clipped track fitting and multi-scale gap detection; found \textbf{13} candidate underdensities, the strongest at \textbf{7.9$\sigma$} (Indus) and \textbf{7.5$\sigma$} (C-12).
\item Fitted stream ages by \textbf{length matching} (Fardal et al.\ 2015 particle spray in a Milky Way potential, following Lu et al.\ 2026), reproducing four published ages to within \textbf{0.02--0.29~Gyr}; derived new ages of \textbf{1.18~Gyr} (\textbf{0.54--2.58}) for Ylgr and \textbf{0.13~Gyr} (\textbf{0.07--0.22}) for C-12 over a progenitor-mass grid.
\item Computed per-stream \textbf{subhalo encounter rates} across \textbf{103} TNG50-1 Milky Way analogues from orbit-derived shells, fitted ages, and per-subhalo relative speeds; Orphan-Chenab has the highest median rate (\textbf{0.032}) with encounters in \textbf{98\%} of halos, versus \textbf{6\%} for C-12 and \textbf{1\%} for NGC~6397.
\end{itemize}
 
\vspace{0.4cm}
 
\textbf{Winter School: MCMC Methods}
\hfill Jan -- Mar 2026 \\
Indian Institute of Technology Roorkee (IITR) \\
Coordinator: Sumedh Gosavi (IITR)
 
\begin{itemize}
\item Implemented \textbf{Metropolis}, \textbf{Metropolis--Hastings} (with asymmetric proposals) and \textbf{Gibbs} samplers, alongside accept--reject and importance sampling, and studied the theory of Monte Carlo and Markov Chain Monte Carlo methods.
\item Applied MCMC to Bayesian parameter estimation (e.g., coin-bias inference with credible intervals), examining posterior distributions and convergence using trace plots, burn-in and acceptance rates.
\item Simulated the 2D Ising model across temperatures, and built a Metropolis-based 1D traffic-jam simulator as the final project, showing how a single slow car triggers backward-propagating congestion.
\end{itemize}
 
\vspace{0.4cm}
 
\textbf{Capstone Project: Exoplanet Habitability Analysis}
\hfill May -- Jul 2025 \\
BRICS Astronomy 2025 \\
Coordinator: Sally Macfarlane, Inter-University Institute for Data Intensive Astronomy (IDIA)
 
\begin{itemize}
\item Cleaned exoplanet datasets and performed feature selection and exploratory analysis using \textbf{NumPy}, \textbf{pandas}, \textbf{Matplotlib}, and \textbf{Seaborn} to study indicators of planetary habitability.
\item Trained and evaluated \textbf{scikit-learn} classification models for habitability prediction.
\end{itemize}
 
% --------- Skills ----------
\section*{Skills}
 
\textbf{Programming:} Python \\
\textbf{Libraries:} NumPy, pandas, SciPy, Matplotlib, Seaborn, scikit-learn, Astropy, h5py, gala, galstreams \\
\textbf{Data and Simulations:} Gaia DR3 (catalogue querying, coordinate transformations), IllustrisTNG \\
\textbf{Methods:} Statistical inference, MCMC, KDE-based density estimation, stellar-stream member selection, machine learning \\
\textbf{Tools:} Git, LaTeX
 
% --------- Achievements ----------
\section*{Achievements}
 
\begin{itemize}
\item Awarded the \textbf{INSPIRE Scholarship} for academic excellence.
\item Ranked in the \textbf{top 1\%} in the ISC (Class XII) examinations.
\end{itemize}
 
\end{document}
 
